\section{Perulangan: for, while, do-while, foreach, break, dan continue}

Perulangan (\textit{loop}) adalah struktur kontrol yang memungkinkan eksekusi blok kode secara berulang selama kondisi tertentu terpenuhi. Tanpa perulangan, tugas seperti menampilkan 100 baris data atau memproses setiap elemen array harus ditulis sebagai 100 baris perintah terpisah. PHP menyediakan empat jenis perulangan: \code{for}, \code{while}, \code{do-while}, dan \code{foreach}, masing-masing cocok untuk skenario yang berbeda \cite{phpManual}.

Perulangan \code{for} paling cocok ketika jumlah iterasi sudah diketahui sebelumnya. Sintaksnya terdiri dari tiga bagian: inisialisasi, kondisi, dan increment/decrement. Perulangan \code{while} memeriksa kondisi di awal setiap iterasi; cocok ketika jumlah iterasi tidak diketahui. Perulangan \code{do-while} mirip dengan \code{while}, tetapi kondisi diperiksa di akhir, sehingga blok kode dieksekusi \textbf{minimal satu kali} meskipun kondisi awalnya sudah \code{false} \cite{w3schoolsPhp}.

\begin{figure}[ht]
\centering
\begin{tikzpicture}[
  node distance=1cm,
  decision/.style={diamond, draw, fill=orange!20, text width=2cm, text badly centered, inner sep=1pt, aspect=2, font=\small},
  block/.style={rectangle, draw, fill=blue!15, rounded corners, text width=2.5cm, minimum height=0.7cm, text centered, font=\small},
  line/.style={draw, -Stealth, thick},
  startstop/.style={rectangle, draw, fill=green!15, rounded corners, minimum width=1.5cm, minimum height=0.6cm, text centered, font=\small},
  title/.style={font=\small\bfseries}
]
  % FOR loop
  \node[title] (tfor) {for};
  \node[block, below=0.5cm of tfor] (init) {Inisialisasi};
  \node[decision, below=of init] (cond) {Kondisi?};
  \node[block, below=of cond] (body) {Eksekusi blok};
  \node[block, right=1.5cm of body] (incr) {Increment};

  \path[line] (init) -- (cond);
  \path[line] (cond) -- node[right, font=\scriptsize]{Ya} (body);
  \path[line] (body) -- (incr);
  \path[line] (incr) |- (cond);

  % WHILE loop
  \node[title, right=5cm of tfor] (twhile) {while};
  \node[decision, below=0.8cm of twhile] (wcond) {Kondisi?};
  \node[block, below=of wcond] (wbody) {Eksekusi blok};

  \path[line] (wcond) -- node[right, font=\scriptsize]{Ya} (wbody);
  \path[line] (wbody.west) -- ++(-0.8,0) |- (wcond.west);
\end{tikzpicture}
\caption{Flowchart perulangan for dan while}
\end{figure}

\begin{webcode}[language=PHP, caption={Perulangan for, while, dan do-while}]
<?php
// for: jumlah iterasi diketahui
echo "Perulangan for:\n";
for ($i = 1; $i <= 5; $i++) {
    echo "Iterasi ke-$i\n";
}

// while: kondisi diperiksa di awal
echo "\nPerulangan while:\n";
$j = 1;
while ($j <= 3) {
    echo "j = $j\n";
    $j++;
}

// do-while: minimal 1 kali eksekusi
echo "\nPerulangan do-while:\n";
$k = 10;
do {
    echo "k = $k\n";  // dieksekusi 1 kali
    $k++;
} while ($k < 5);     // kondisi false, tapi sudah dieksekusi
?>
\end{webcode}

Perulangan \code{foreach} dirancang khusus untuk mengiterasi array dan merupakan cara paling idiomatik untuk memproses setiap elemen array di PHP. Sintaks \code{foreach (\$array as \$nilai)} mengambil nilai saja, sedangkan \code{foreach (\$array as \$kunci => \$nilai)} mengambil kunci dan nilai. Perulangan ini lebih aman daripada \code{for} karena tidak memerlukan pengelolaan indeks secara manual dan tidak risiko \textit{index out of bounds} \cite{phpManual}.

\begin{webcode}[language=PHP, caption={Perulangan foreach untuk array}]
<?php
// Array terindeks
$buah = ["Apel", "Mangga", "Jeruk"];
foreach ($buah as $item) {
    echo "Buah: $item\n";
}

// Array asosiatif
$mahasiswa = ["nama" => "Rizki", "nim" => "2023001", "ipk" => 3.5];
foreach ($mahasiswa as $kunci => $nilai) {
    echo "$kunci: $nilai\n";
}
?>
\end{webcode}

Kata kunci \code{break} dan \code{continue} memberikan kontrol tambahan di dalam perulangan. \code{break} menghentikan seluruh perulangan dan langsung keluar dari loop. \code{continue} melompati sisa kode di iterasi saat ini dan langsung menuju iterasi berikutnya. Keduanya sangat berguna untuk mengoptimalkan logika loop, misalnya menghentikan pencarian setelah item ditemukan atau melewati elemen yang tidak memenuhi kriteria \cite{phpRightWay}.

\begin{catatan}
\textbf{Hindari Infinite Loop:} Pastikan kondisi perulangan pasti akan bernilai \code{false} pada suatu saat. Jika lupa menginkrementasi counter atau kondisi tidak pernah terpenuhi, program akan berjalan tanpa henti dan server PHP akan timeout. Gunakan \code{set\_time\_limit()} untuk membatasi waktu eksekusi saat debugging.
\end{catatan}

